\scanpage{28}
Alors il existe un unique $\xi\in M$ tel que l'on ait
\[
\xi_k=c_{\deleted{\unsure{k}}i,k-i}\xi\qquad\left(=\binom{k}{i}\xi\right)\qquad\text{pour tt }k\geq N\tag*{(2)}
\]
(\uline{NB} $c_{k,l}$ désigne le coefficient binômial $\frac{(k+l)!}{k!\,l!}$)

\uline{Dém} \textcircled{1} Réduction au cas où $\exists p$ nb premier, tel que $\deleted{\unsure{p}}n\,\mathrm{id}_M$ bijectif $\forall n\in\mathbb Z$ tel que $(n,p)=1$, i.e. $M$ un $\mathbb Z_p$ module ($\mathbb Z_p\overset{\mathrm{def}}{=}\mathbb Z_{(p\mathbb Z)}$).

Supposons que ce soit prouvé dans ce cas, et qu'on sache qu'il existe alors $k=k(p)$ tel que $\deleted{\unsure{\xi}}c_{\deleted{\unsure{k}}i,k-i}\not\equiv0\ (p)$ (i.e $c_{k,k-i}\in\mathbb Z_p^{\unsure{\times|*}}$]. Prouvons le th. en général. Les $c_{k,k-i}$ ($k\geq N$) ne sont pas divisibles par un $=$ nb premier, donc sont premiers dans leur ensemble, donc $\exists k_1\ldots k_r\geq N$ et des entiers $m_1,\ldots m_r\in\mathbb Z$ tels que
\[
\sum_{\alpha=1}^{r}m_\alpha c_{k_\alpha,k_\alpha-i}=\unsure{1|d}.
\]
Donc les relations (2) impliquent, en multipliant par les $m_\alpha$ et ajoutant
\[
\xi=\sum_1^r m_\alpha\xi_{k_\alpha}
\]
d'où l'unicité de $\xi$. Prouvons qu'avec ce choix de $\xi$, on a bien (2) i.e.
\[
\xi_k=c_{k,k-i}\sum_1^{\unsure{r}} m_\alpha\xi_{k_\alpha}
\]
Pour prouver que cette relation est vraie, \deleted{\unsure{p|k ?}} il suffit de prouver qu'elle est vraie dans chaque $M_p=M\otimes_{\mathbb Z}\mathbb Z_p$ ($p\in\mathbb P$). Mais c'est une conséquence immédiate du fait que le th. est vrai pour $M_p$.
% archive sheet number: 27

